% Corrected French for printed p.129.
% Corrections: missing une before conséquence; utile -> utiles.
\SourcePage{129}

\LayerTarget{tohoku:section:1.5:p1}
\noindent dans $C$, est \textit{exact à droite}. Il est même exact si $I$ est
fini, mais pas nécessairement si $I$ est infini, car la somme directe d'une
famille infinie de monomorphismes n'est pas nécessairement un monomorphisme,
comme nous l'avons remarqué au N\textsuperscript{o}1 (pour la situation duale).
D'où l'axiome suivant :

\LayerTarget{tohoku:section:1.5:AB4}
\noindent\textbf{AB 4)} \textit{L'axiome AB 3) est vérifié, et la somme directe
d'une famille de monomorphismes est un monomorphisme.}

L'axiome suivant est strictement plus fort que AB 4) :

\LayerTarget{tohoku:section:1.5:AB5}
\noindent\textbf{AB 5)} \textit{L'axiome AB 3) est vérifié, et si
$(A_i)_{i\in I}$ est une famille filtrante croissante de sous-trucs d'un
$A\in C$, $B$ un sous-truc quelconque de $A$, on a}

\[
  \left(\sum_i A_i\right)\cap B=\sum_i(A_i\cap B).
\]

\noindent (On a noté, conformément aux usages, $\sum A_i$ le sup des $A_i$, et
$P\cap Q$ le inf des sous-trucs $P$ et $Q$ de $A$).

\LayerTarget{tohoku:section:1.5:AB5:equivalent}
AB 5) peut encore s'exprimer ainsi : \textit{AB 3) est satisfait, et si
$A\in C$ est le sup d'une famille filtrante croissante de sous-trucs $A_i$, et
si pour tout $i$ on se donne un morphisme
$u_i:A_i\longrightarrow B$ tel que pour $A_i\subset A_j$, $u_i$ soit induit
par $u_j$, alors il existe un morphisme $u$} (évidemment unique) \textit{de $A$
dans $B$ qui induise les $u_i$}. Notons enfin l'axiome suivant qui renforce
encore AB 5), mais dont nous n'aurons pas à nous servir dans ce travail :

\LayerTarget{tohoku:section:1.5:AB6}
\noindent\textbf{AB 6)} \textit{L'axiome AB 3) est vérifié, et pour tout
$A\in C$ et toute famille $(B^j)_{j\in J}$ de familles
$B^j=(B_i^j)_{i\in I_j}$ filtrantes croissantes de sous-trucs $B_i^j$ de $A$,
on a}

\[
 \bigcap_{j\in J}\left(\sum_{i\in I_j}B_i^j\right)
 =\sum_{(i_j)\in\prod_{j\in J}I_j}
   \left(\bigcap_{j\in J}B_{i_j}^j\right).
\]

\noindent (Cet axiome inclut implicitement l'existence du inf d'une famille
quelconque de sous-trucs de $A$).

\LayerTarget{tohoku:section:1.5:p2}
Nous laissons au lecteur le soin d'énoncer les axiomes duals AB 3*), AB 4*),
AB 5*) et AB 6*), relatifs aux produits infinis, des axiomes précédents.
Signalons à titre d'exemple que la catégorie des groupes abéliens (ou plus
généralement, la catégorie des modules sur un anneau avec unité fixé)
satisfait, relativement aux sommes directes, à l'axiome le plus fort AB 6),
elle satisfait de plus aux axiomes AB 3*) et AB 4*), mais non à AB 5*). Les
faits sont inverses pour la catégorie duale, qui par la dualité de Pontrjagin
est isomorphe à la catégorie des groupes abéliens topologiques compacts. (Cela
montre que AB 5*) n'est pas une conséquence de AB 4*), donc AB 5) n'est pas non
plus une conséquence de AB 4)). La catégorie abélienne des faisceaux de groupes
abéliens sur un espace topologique fixé $X$ satisfait à l'axiome AB 5) et
AB 3*), mais non à AB 4*), car nous avons déjà remarqué qu'un produit de
morphismes surjectifs n'est pas nécessairement un morphisme surjectif.
Remarquons pour finir que si $C$ est une catégorie satisfaisant à la fois à
AB 5) et AB 5*), alors $C$ est réduit aux éléments nuls (car on voit alors
facilement que pour $A\in C$, le morphisme canonique
$A^{(I)}\longrightarrow A^I$ est un isomorphisme, et on vérifie que ceci n'est
possible que si $A$ est nul).

\LayerTarget{tohoku:section:1.5:p3}
Les axiomes précédents nous seront surtout utiles pour l'étude des limites
inductives et projectives dont nous aurons besoin pour donner des conditions

\TohokuContinuation{130}
